Consideriamo il problema di fusione a livello di famiglia CWE. Per ogni artefatto di codice $x$ e per una famiglia fissata $f$, ciascun analizzatore statico $T_i$ produce una decisione binaria
\begin{equation}
e_i^f(x) \in \{0,1\},
\end{equation}
dove $e_i^f(x)=1$ indica che $T_i$ segnala la famiglia $f$ su $x$, mentre $e_i^f(x)=0$ indica silenzio rispetto a quella famiglia. Con $L$ strumenti, l'output congiunto e' il vettore di comportamento
\begin{equation}
E_f(x)=\bigl(e_1^f(x),e_2^f(x),\ldots,e_L^f(x)\bigr).
\end{equation}
La variabile di verita' e' $y_f(x)\in\{0,1\}$: $y_f(x)=1$ significa che l'artefatto e' realmente vulnerabile per $f$. Nel seguito useremo anche le classi $\Vuln$ e $\Safe$, rispettivamente ``vulnerabile'' e ``safe''. Ogni tecnica produce uno score $s_f(x)\in[0,1]$ interpretabile come evidenza a favore di $\Vuln$; la decisione finale e'
\begin{equation}
\hat{y}_f(x)=
\begin{cases}
1, & s_f(x)\geq \tau,\\
0, & s_f(x)<\tau,
\end{cases}
\end{equation}
con soglia di default $\tau=0.5$.

La calibrazione e' stimata per coppia strumento--famiglia. Dal punto di vista della repository, questa informazione e' prodotta da \texttt{analysis/calibration.py}: per ogni $(T_i,f)$ si calcolano la matrice di confusione one-vs-rest e metriche come precisione positiva (\texttt{ppv}), valore predittivo negativo (\texttt{npv}), tasso di falsi positivi (\texttt{fpr}), tasso di falsi negativi (\texttt{fnr}), sensibilita' e specificita'. Una coppia $(T_i,f)$ e' considerata supportata se lo strumento ha almeno una segnalazione per quella famiglia nel set di calibrazione; le strategie pesate usano solo strumenti supportati.

Useremo lo stesso esempio numerico in tutte le sezioni. Fissiamo una famiglia $f$ e tre analizzatori: CodeQL ($T_1$), Semgrep ($T_2$) e Joern ($T_3$). Sul frammento $x^\star$ osserviamo
\begin{equation}
E_f(x^\star)=(1,0,0),
\end{equation}
cioe' CodeQL segnala la CWE, mentre Semgrep e Joern restano silenti. Nel set di calibrazione assumiamo una prevalenza
\begin{equation}
P(\Vuln)=0.40,\qquad P(\Safe)=0.60.
\end{equation}
Le affidabilita' usate nell'esempio sono riportate in Tabella~\ref{tab:shared_example}.

\begin{table}[h]
\centering
\small
\begin{tabular}{lccccccc}
\toprule
Strumento & $e_i$ & $P(1\mid\Vuln)$ & $P(1\mid\Safe)$ & $P(0\mid\Vuln)$ & $P(0\mid\Safe)$ & \texttt{ppv} & \texttt{npv}\\
\midrule
CodeQL  & 1 & 0.80 & 0.15 & 0.20 & 0.85 & 0.78 & 0.86\\
Semgrep & 0 & 0.70 & 0.15 & 0.30 & 0.85 & 0.76 & 0.81\\
Joern   & 0 & 0.60 & 0.10 & 0.40 & 0.90 & 0.80 & 0.77\\
\bottomrule
\end{tabular}
\caption{Esempio condiviso per una famiglia CWE. Lo stesso vettore $E=(1,0,0)$ viene valutato da tutte le strategie.}
\label{tab:shared_example}
\end{table}

Per le tecniche che apprendono dai pattern completi useremo anche due informazioni aggiuntive. Primo, per Logistic Regression assumiamo coefficienti illustrativi
\begin{equation}
w_0=-1.2,\qquad w_1=2.1,\qquad w_2=0.8,\qquad w_3=0.5.
\end{equation}
Secondo, per BKS assumiamo che il pattern $E=(1,0,0)$ sia comparso $50$ volte nel set di calibrazione, con $42$ esempi vulnerabili e $8$ safe. Questi numeri non descrivono una specifica run sperimentale: servono a rendere confrontabile, sezione per sezione, l'effetto delle diverse ipotesi di fusione.
